% from ega4-18-fr.tex
\oldpage[IV-4]{114}

Cela étant, pour que l'homomorphisme $u$ soit bijectif, il faut et il suffit que $\bigwedge^n u$ soit bijectif, ou encore que la section $d_{X/Y}$ ait un germe \emph{inversible} en tout point $y\in Y$, ce qui s'écrit aussi $d_{X/Y}(y)\neq 0$ pour tout $y$ ($0_{\mathrm I}$, 5.5.2). Il revient aussi au même de dire que l'Idéal discriminant $\mathscr{D}_{X/Y}$ est \emph{égal à $\sh{O}_Y$}.

La terminologie de « revêtement » introduite dans (18.2.7, (i)), se justifie par la proposition suivante:

\begin{proposition}[18.2.8]
\label{IV.18.2.8}
Soit $f:X\to Y$ un morphisme étale, séparé et de type fini et pour tout $y\in Y$, soit $n(y)$ le nombre géométrique des points de $f^{-1}(y)$. Alors la fonction $y\mapsto n(y)$ est semi-continue inférieurement dans $Y$. Pour qu'elle soit continue en un point $y$ (donc constante dans un voisinage de $y$), il faut et il suffit qu'il existe un voisinage ouvert $U$ de $y$ tel que la restriction $f^{-1}(U)\to U$ de $f$ soit un morphisme fini (étale).
\end{proposition}

Comme $f$ est quasi-fini (17.6.1) et localement de présentation finie, il revient au même de dire que $f$ est fini ou que $f$ est propre (8.11.1): en outre, chaque fibre $f^{-1}(y)$ est géométriquement réduite sur $\kappa(y)$. Les conclusions résultent donc de (15.5.9, (i) et (ii)) et de ce que $f$ est plat.

\begin{corollary}[18.2.9]
\label{IV.18.2.9}
Soient $Y$ un préschéma connexe, $f:X\to Y$ un morphisme étale, séparé et de type fini. Pour que $f$ soit fini (autrement dit, que $X$ soit un revêtement étale de $Y$), il faut et il suffit que toutes les fibres de $f$ aient même nombre géométrique de points.
\end{corollary}

\begin{remarks}[18.2.10]
\label{IV.18.2.10}
\begin{enumerate}
  \item[(i)] L'exemple de la « droite affine avec un point dédoublé » (I, 5.5.11) montre qu'un morphisme étale de type fini de préschémas noethériens peut être \emph{non séparé}; pour cet exemple, la première assertion de (18.2.8) n'est plus exacte.

  \item[(ii)] Pour qu'un morphisme séparé, de type fini et étale $f:X\to Y$ fasse de $X$ un revêtement \emph{localement trivial}, il faut et il suffit que pour tout $x\in X$ il existe un voisinage ouvert $U$ de $y=f(x)$ et une $U$-section $g$ de $f^{-1}(U)$ telle que $g(y)=x$. En effet, la condition est évidemment nécessaire; le fait qu'elle est suffisante résulte de ce que toute fibre $f^{-1}(y)$ est finie (17.6.1) et (I, 6.2.2), de la caractérisation des sections d'un $Y$-schéma étale (17.9.3) et de la prop. (18.2.8).
\end{enumerate}
\end{remarks}

\subsection{Algèbres finies et étales.}
\EGAReaderTarget{ega4.c04.subsec.18-3}{subsection:IV.18.3}

\begin{proposition}[18.3.1]
\label{IV.18.3.1}
Soient $A$ un anneau, $B$ une $A$-algèbre de présentation finie.
\begin{enumerate}
  \item[(i)] Pour que $B$ soit une $A$-algèbre non ramifiée, il faut et il suffit que $B$ soit un $A$-module de présentation finie et que $B$ soit un $(B\otimes_A B)$-module projectif.
  \item[(ii)] Supposons de plus que $B$ soit une $A$-algèbre finie. Pour que $B$ soit une $A$-algèbre étale, il faut et il suffit que $B$ soit un $A$-module projectif et un $(B\otimes_A B)$-module projectif.
\end{enumerate}
\end{proposition}

Bien entendu, la structure de $(B\otimes_A B)$-module sur $B$ est celle provenant de la structure de $(B\otimes_A B)$-algèbre sur $B$ correspondant au $A$-homomorphisme canonique d'anneaux $B\otimes_A B\to B$, qui est surjectif et de noyau $\mathfrak J_{B/A}$ (0, 20.4.1).

(i) Dire que le morphisme $\Spec(B)\to\Spec(A)$ est localement de présentation finie équivaut à dire que $B$ est une $A$-algèbre de présentation finie (1.4.6). Dire

